\documentclass[12pt,reqno]{amsart}
\usepackage{../handout, ../customcommands}
\usepackage{caption}

\hdr{MAT 240}{\S\S5.3-5.4 Hypergeometric and Poisson distributions}

\begin{document}

\textbf{Suggested problems:} \S5.3: 2-5, \S5.4: 2, 3, 6, 7

\bigskip
\hrule

\bigskip
\prob Suppose that a box contains $A$ red balls and $B$ blue balls. Suppose that $n$ balls are selected from the box (without replacement) and let $X$ be the number of red balls obtained. What is the probability mass function of $X$?

\vfill
\prob Over a span of $t$ hours, let $X$ denote the number of emails that arrive in your inbox. Assuming that the emails arrive at random and independently of each other, and that the mean rate at which they arrive is $\lambda$ (measured in units of reciprocal hours, or hour$^{-1}$), what is the distribution of the random variable $X$?

\vfill
\newpage
\prob Suppose that $X\sim \mathcal{P}ois(\mu)$. Compute the moment-generating function $\psi(t)$ of $X$, as well as $E(X)$ and $V(X)$.

\vfill
\prob A store owner believes that customers arrive at his store at a rate of 4.5 customers per hour on average.

\medskip
\begin{enumerate}
\item Describe a distribution that models the number $X$ of customers that arrive at the store during a two-hour time interval.\vfill
\item Compute the probabilities that one customer arrives in the time interval, and that ten customers arrive.\vfill
\item Compute the probabilities that at most five customers arrive in the time interval, and that more than ten arrive.\vfill
\end{enumerate}


\end{document}